\exercise{14.5}{Linearity of series and failure of a termwise product formula}
\goal If $\sum a_n=A$ and $\sum b_n=B$ are finite real sums,
prove the two linearity formulas and examine the proposed product
formula.

\step{Preparation: use partial sums, whose limits are given.}
Write $S_N=\sum_{n=1}^{N}a_n$ and $T_N=\sum_{n=1}^{N}b_n$.
The hypotheses mean $S_N\to A$ and $T_N\to B$. These are finite
limits, so the sum and constant multiple laws from Section~9 apply.

\step{(a) Addition of two series.}
Finite distributivity gives
\[
 \sum_{n=1}^{N}(a_n+b_n)=S_N+T_N.
\]
Taking limits yields $A+B$. By the definition of a series sum,
this proves $\sum_{n=1}^{\infty}(a_n+b_n)=A+B$.

\step{(b) Multiplication by a fixed real constant.}
For any $k\in\R$ and every $N$,
\[
 \sum_{n=1}^{N}ka_n=kS_N.
\]
The right side tends to $kA$, including when $k=0$ or $k<0$.
Hence $\sum_{n=1}^{\infty}ka_n=kA$.

\step{(c) Test the termwise product formula on finite series.}
Let $a_1=1$ and $a_n=0$ for $n\ge2$. Let $b_2=1$ and
$b_n=0$ for every $n\ne2$. Each series has sum $1$, so $AB=1$.
At $n=1$ the second factor is zero; at $n=2$ the first factor is
zero; at every other index both are zero. Thus $a_nb_n=0$ for
all $n$, and
\[
 \sum_{n=1}^{\infty}a_nb_n=0\ne1=AB.
\]
The conjectured identity is false even for series with only one
nonzero term each. The finite product $S_NT_N$ contains all the
terms $a_i b_j$ with $1\le i,j\le N$, whereas the partial sum
$\sum_{n=1}^{N}a_nb_n$ includes only the terms with equal indices.
The product limit law applies to $S_NT_N$, not to that shorter sum.
\answer{The two linearity identities hold. The termwise product
identity does not follow and is false in general.}

\exercise{14.6}{Absolutely summable terms times a bounded sequence}
\goal Prove (a) that $\sum a_nb_n$ converges when
$\sum\abs{a_n}$ converges and $(b_n)$ is bounded, and (b) recover
Corollary~14.7.

\step{(a), first step: put one bound on every second factor.}
Boundedness supplies a finite $M>0$ with $\abs{b_n}\le M$ for
every $n$. We may insist on $M>0$ by enlarging a given bound,
even if every $b_n$ is zero.

\step{(a), second step: verify the Cauchy criterion for the product series.}
Let $\varepsilon>0$. Since the nonnegative series
$\sum\abs{a_n}$ converges, its Cauchy criterion gives an index
$N$ such that
\[
 n\ge m>N\quad\Longrightarrow\quad
 \sum_{j=m}^{n}\abs{a_j}<\frac{\varepsilon}{M}.
\]
For any such $m,n$, the finite triangle inequality and the bound
on $b_j$ imply
\[
 \abs{\sum_{j=m}^{n}a_jb_j}
 \le\sum_{j=m}^{n}\abs{a_j}\abs{b_j}
 \le M\sum_{j=m}^{n}\abs{a_j}
 <M\frac{\varepsilon}{M}=\varepsilon.
\]
This is exactly the Cauchy criterion for $\sum a_nb_n$.
Theorem~14.4 therefore proves convergence. The same estimate
without the outer absolute value also proves the Cauchy criterion
for $\sum\abs{a_nb_n}$, so the product series actually converges
absolutely.

\step{(b) Use a constant second sequence.}
Choose $b_n=1$ for all $n$. This sequence is bounded and
$a_nb_n=a_n$. Part (a) then says that convergence of
$\sum\abs{a_n}$ implies convergence of $\sum a_n$, precisely
Corollary~14.7.
\answer{Every bounded multiplier preserves absolute summability;
the constant multiplier $1$ gives the cited corollary.}

\exercise{14.7}{Powers greater than one of a summable nonnegative sequence}
\goal If $a_n\ge0$, $\sum a_n$ converges, and $p>1$, prove that
$\sum a_n^p$ converges.

\step{1. A convergent series has terms tending to zero.}
Corollary~14.5 gives $a_n\to0$. Hence there is $N_0$ such that
$0\le a_n<1$ for all $n>N_0$.

\step{2. Compare the powers only where that comparison is justified.}
If $a_n=0$, then $a_n^p=0\le a_n$. If $0<a_n\le1$, the
positive exponent $p-1$ gives $a_n^{p-1}\le1$. Multiplication
by $a_n>0$ yields
\[
 0\le a_n^p=a_na_n^{p-1}\le a_n\qquad(n>N_0).
\]
The restriction to sufficiently late indices matters: a number
larger than $1$ has its size increased by taking a power $p>1$.

\step{3. Include the finite initial part.}
Let $A=\sum_{n=1}^{\infty}a_n<+\infty$ and
$C=\sum_{n=1}^{N_0}a_n^p$, both nonnegative and finite.
For $N>N_0$,
\[
 \sum_{n=1}^{N}a_n^p
 =C+\sum_{n=N_0+1}^{N}a_n^p
 \le C+\sum_{n=N_0+1}^{N}a_n\le C+A.
\]
For $N\le N_0$, the partial sum is at most $C\le C+A$ as well.
The partial sums of $\sum a_n^p$ are increasing because their
terms are nonnegative, and the displayed inequalities bound them
above. The monotone sequence theorem proves a finite limit.
\answer{The series $\sum a_n^p$ converges for every real $p>1$.}

\exercise{14.8}{A geometric-mean comparison for two nonnegative series}
\goal If $a_n,b_n\ge0$ and both $\sum a_n$ and $\sum b_n$
converge, prove convergence of $\sum\sqrt{a_nb_n}$.

\step{1. Derive the needed pointwise inequality.}
The square of any real number is nonnegative. Apply this to
$\sqrt{a_n}-\sqrt{b_n}$, whose square expands to
\[
 0\le(\sqrt{a_n}-\sqrt{b_n})^2
   =a_n+b_n-2\sqrt{a_nb_n}.
\]
Rearranging and dividing by $2>0$ gives
\[
 0\le\sqrt{a_nb_n}\le\frac{a_n+b_n}{2}\le a_n+b_n.
\]
The last inequality uses $a_n+b_n\ge0$. The square-root identity
also holds when one or both of the numbers are zero.

\step{2. Apply the Comparison Test to a convergent upper series.}
By Exercise~14.5, $\sum(a_n+b_n)$ converges, with sum equal to
the sum of the two original sums. Its terms are nonnegative.
The displayed pointwise inequality therefore satisfies all the
hypotheses of the Comparison Test, which proves convergence of
$\sum\sqrt{a_nb_n}$.
\answer{The geometric-mean series converges. In fact its sum is
at most one half of the sum of the two given series.}

\exercise{14.9}{Changing finitely many terms}
\goal Suppose $a_n=b_n$ except for finitely many indices. Prove
that $\sum a_n$ and $\sum b_n$ either both converge or both
diverge.

\step{1. Place all exceptional indices before a single cutoff.}
If the exceptional set is nonempty, let $N_0$ be its maximum.
If the set is empty, take $N_0=1$. In either case $a_n=b_n$
for every $n>N_0$. Define the finite real constant
\[
 D=\sum_{n=1}^{N_0}(a_n-b_n).
\]

\step{2. Compare partial sums exactly.}
For each $N\ge N_0$, all the later differences are zero, so
\[
 \sum_{n=1}^{N}a_n-\sum_{n=1}^{N}b_n
 =\sum_{n=1}^{N_0}(a_n-b_n)=D.
\]
Thus the two partial-sum sequences differ by a fixed constant
from that point onward.

\step{3. Prove both directions of convergence.}
If $\sum a_n=A$ is finite, the last identity gives
$\sum_{n=1}^{N}b_n=\sum_{n=1}^{N}a_n-D\to A-D$, so the
second series converges. Conversely, if $\sum b_n=B$ is finite,
then $\sum_{n=1}^{N}a_n=\sum_{n=1}^{N}b_n+D\to B+D$.
Therefore convergence of either series implies convergence of the
other. If one diverges and the other converged, one of these
implications would be contradicted; hence divergence is shared too.
\answer{Finite changes preserve convergence. They can change a
finite sum by the constant $D$, which is why equality of sums is
not asserted.}

\exercise{14.10}{Divergence detected by roots but not by ratios}
\goal Give a series for which the Root Test proves divergence
while neither conclusion of the Ratio Test applies.

\step{1. Define positive terms with alternating adjacent ratios.}
Take
\[
 a_n=2^{\,n+(-1)^n}\qquad(n\ge1).
\]
Every term is positive, so its root and ratio are defined.
The $n$th root is $a_n^{1/n}=2\cdot2^{(-1)^n/n}$.
Since $-1/n\le(-1)^n/n\le1/n$ and powers of $2$ preserve
order,
\[
 2\cdot2^{-1/n}\le a_n^{1/n}\le2\cdot2^{1/n}.
\]
Both outer expressions tend to $2$, using the root limit for the
fixed positive number $2$ and its reciprocal. The squeeze theorem
gives $a_n^{1/n}\to2$.

\step{2. Apply the divergent part of the Root Test.}
The root limit superior is $2>1$, so the Root Test proves
divergence of $\sum a_n$.

\step{3. Check both possible ratio-test inequalities.}
When $n$ is even, $(-1)^n=1$ and $(-1)^{n+1}=-1$, giving
\[
 \frac{a_{n+1}}{a_n}
 =2^{(n+1)-1-(n+1)}=2^{-1}=\frac12.
\]
When $n$ is odd, the two signs are reversed, giving
\[
 \frac{a_{n+1}}{a_n}
 =2^{(n+1)+1-(n-1)}=2^3=8.
\]
Each value occurs infinitely often, and there are no other ratio
values. Hence the ratio limit inferior is $1/2$ and its limit
superior is $8$. The convergence condition, ratio limit superior
less than $1$, fails. The divergence condition, ratio limit
inferior greater than $1$, fails as well. Thus the Ratio Test is
inconclusive even though the Root Test settles the question.
\answer{The series $\sum 2^{n+(-1)^n}$ has the required properties:
root limit $2$, ratio limit inferior $1/2$, and ratio limit superior $8$.}

\exercise{14.11}{A constant successive ratio gives a geometric series}
\goal Suppose every $a_n\ne0$ and $a_{n+1}/a_n=r$ is independent
of $n$. Prove that $\sum a_n$ is geometric.

\step{1. Convert the quotient statement to a recurrence.}
Since $a_n\ne0$, multiplication by $a_n$ gives
$a_{n+1}=ra_n$ for every $n$. Also $r=a_2/a_1\ne0$ because
both $a_1$ and $a_2$ are nonzero.

\step{2. Solve the recurrence by induction.}
At $n=1$, the formula $a_n=a_1r^{n-1}$ reads $a_1=a_1r^0$,
which holds. Suppose $a_n=a_1r^{n-1}$ for some $n\ge1$.
Then
\[
 a_{n+1}=ra_n=r(a_1r^{n-1})=a_1r^n.
\]
This is the same formula with $n$ replaced by $n+1$, so induction
proves it for every positive integer.

\step{3. Recognize the series and its possible sum.}
The series is therefore
\[
 \sum_{n=1}^{\infty}a_n
 =\sum_{n=1}^{\infty}a_1r^{n-1}
 =a_1+a_1r+a_1r^2+\cdots,
\]
which is the definition of a geometric series with first term
$a_1$ and common ratio $r$. It converges when $\abs r<1$ and
then has sum $a_1/(1-r)$. If $\abs r\ge1$, the nonzero first
term ensures that the terms do not tend to zero, so it diverges.
\answer{A constant successive ratio forces the exact geometric
form $a_n=a_1r^{n-1}$.}

\exercise{14.12}{Selecting an absolutely summable subsequence}
\goal From $\liminf\abs{a_n}=0$, construct strictly increasing
indices $n_k$ for which $\sum_{k=1}^{\infty}a_{n_k}$ converges.

\step{1. Arbitrarily small absolute values occur after every finite prefix.}
Fix a number $\eta>0$ and an index cutoff $N\ge0$. If no
$n>N$ satisfied $\abs{a_n}<\eta$, every term beyond $N$
would have $\abs{a_n}\ge\eta$. All later tail infima would then
be at least $\eta$, which would imply
$\liminf\abs{a_n}\ge\eta>0$, a contradiction. Thus such an
index always exists.

\step{2. Select a geometric error bound at each stage.}
Put $n_0=0$. After $n_{k-1}$ has been chosen, use Step~1 with
$N=n_{k-1}$ and $\eta=2^{-k}$ to choose
\[
 n_k>n_{k-1},\qquad\abs{a_{n_k}}<2^{-k}.
\]
For definiteness, choose the least eligible index. Existence is
provided by Step~1, and the least positive integer principle makes
the recursive choice well defined. Strict increase of the indices
ensures that $(a_{n_k})$ is a subsequence of the original sequence.

\step{3. Compare the absolute-value series to a geometric series.}
The geometric series satisfies $\sum_{k=1}^{\infty}2^{-k}=1$.
Its terms dominate $\abs{a_{n_k}}\ge0$, so the Comparison Test
gives convergence of $\sum\abs{a_{n_k}}$. Absolute convergence
implies convergence of $\sum a_{n_k}$ by Corollary~14.7.
The geometric bound is essential to this argument: merely choosing
a subsequence whose terms tend to zero would not by itself prove
that its series converges.
\answer{There is an absolutely convergent subsequence series,
obtained by imposing $\abs{a_{n_k}}<2^{-k}$ at stage $k$.}

\exercise{14.13}{Four exact sums from geometric and telescoping identities}
\goal Calculate the two geometric sums in part (a), prove the
identities in parts (b) and (c), and deduce the sum in part (d).

\step{(a) The ratios $2/3$ and $-2/3$, starting at $n=1$.}
For any $\abs r<1$, shifting the index of the geometric series gives
\[
 \sum_{n=1}^{\infty}r^n
 =r\sum_{j=0}^{\infty}r^j=\frac{r}{1-r}.
\]
Both ratios here have absolute value $2/3<1$. Substitution yields
\[
 \sum_{n=1}^{\infty}\left(\frac23\right)^n
 =\frac{2/3}{1-2/3}=2,
 \qquad
 \sum_{n=1}^{\infty}\left(-\frac23\right)^n
 =\frac{-2/3}{1+2/3}=-\frac25.
\]
The negative sign in the second first term matters because its
sum begins at exponent $1$, rather than exponent $0$.

\step{(b) The series with term $1/[n(n+1)]$.}
The finite algebraic identity
\[
 \frac1n-\frac1{n+1}
 =\frac{n+1-n}{n(n+1)}=\frac1{n(n+1)}
\]
gives, for every $N\ge1$,
\[
 \begin{aligned}
 \sum_{n=1}^{N}\frac1{n(n+1)}
 &=(1-\tfrac12)+(\tfrac12-\tfrac13)+\cdots
      +(\tfrac1N-\tfrac1{N+1})\\
 &=1-\frac1{N+1}.
 \end{aligned}
\]
Taking the limit $N\to\infty$ gives the finite sum $1$.

\step{(c) The series with term $(n-1)/2^{n+1}$.}
Put the right side of the following identity over a common
denominator to verify it:
\[
 \frac n{2^n}-\frac{n+1}{2^{n+1}}
 =\frac{2n-(n+1)}{2^{n+1}}
 =\frac{n-1}{2^{n+1}}.
\]
The first term of the series, at $n=1$, is zero. Summing the
identity to $N$ causes all intervening terms to cancel:
\[
 \sum_{n=1}^{N}\frac{n-1}{2^{n+1}}
 =\frac12-\frac{N+1}{2^{N+1}}.
\]
To justify the last term's limit explicitly, set
$u_N=(N+1)/2^{N+1}>0$. For $N\ge1$,
\[
 \frac{u_{N+1}}{u_N}=\frac{N+2}{2(N+1)}\le\frac34,
\]
because the last inequality is equivalent to $N\ge1$.
Induction gives $0<u_N\le u_1(3/4)^{N-1}$, whose upper bound
tends to zero. Hence $u_N\to0$ by the squeeze theorem, and
\[
 \sum_{n=1}^{\infty}\frac{n-1}{2^{n+1}}=\frac12.
\]

\step{(d) Deduce the sum with term $n/2^n$.}
For each index,
\[
 \frac n{2^n}
 =2\frac{n-1}{2^{n+1}}+\frac1{2^n}.
\]
The first series on the right converges by part (c), and the
second is geometric with sum $1$. The linearity formulas of
Exercise~14.5 therefore apply and give
\[
 \sum_{n=1}^{\infty}\frac n{2^n}
 =2\left(\frac12\right)+1=2.
\]
\answer{The requested sums are, in order, $2$, $-2/5$, $1$,
$1/2$, and $2$. All telescoping cancellations were made in finite
partial sums before taking a limit.}

\exercise{14.14}{A blockwise lower comparison for the harmonic series}
\goal Prove divergence of $\sum_{n=1}^{\infty}1/n$ by comparing
it with the displayed repeated-block sequence.

\step{1. Express the displayed sequence by exact index ranges.}
For each integer $j\ge1$, the block
\[
 2^{j-1}\le n\le2^j-1
\]
contains $2^j-1-2^{j-1}+1=2^{j-1}$ indices. On that block
define $a_n=2^{-j}$. Thus the first block is the single term
$1/2$; the second contains two terms $1/4$; the third contains
four terms $1/8$; and the fourth contains eight terms $1/16$.
These blocks exactly reproduce the sequence in the question.

\step{2. Every complete block contributes the same positive amount.}
For the block labeled $j$, its sum is
\[
 2^{j-1}\cdot2^{-j}=\frac12.
\]
The comparison series in the exercise begins at $n=2$, so it
omits the first block and retains all blocks with $j\ge2$.
At the end of block $J\ge2$ its partial sum is
\[
 \sum_{n=2}^{2^J-1}a_n
 =\sum_{j=2}^{J}\frac12=\frac{J-1}{2}.
\]
These values tend to $+\infty$ as $J$ grows. All intermediate
partial sums are increasing, since $a_n>0$. Therefore the entire
comparison series, not only these block-end partial sums,
diverges to $+\infty$.

\step{3. Compare each harmonic term with its block value.}
Given $n\ge2$, there is a unique block index $j\ge2$ with
$2^{j-1}\le n\le2^j-1$. Since $n<2^j$ and both are positive,
\[
 \frac1n>\frac1{2^j}=a_n.
\]
The Comparison Test now shows that $\sum_{n=2}^{\infty}1/n$
diverges to $+\infty$. Adding the first harmonic term, which is
$1$, preserves this divergence.
\answer{The harmonic series diverges to $+\infty$ because it
dominates a positive series whose every complete block adds $1/2$.}
